\section{Limitations}
\label{sec:limitations}

% TRACE: S8-C01 | configs/data_governance.yaml | access_policy, public_release_status; answer/data_provenance/source_registry.json | $.sources; docs/external_dataset_access_and_license_status.md | Decision
The primary limitation is the data-access and provenance boundary.  The fixed corpus combines 10,000 retained local social-media export records with 2,000 author-created, context-augmented derivatives based on VIVID idiom/proverb materials.  Raw and processed source text remains private and is not redistributed; source authorization and license documentation are pending project action.  The paper reports aggregate artifacts and provenance rather than public raw text, and it cannot offer an open-data claim or a public qualitative-example set.  The AIVIVN artifact is a hash-identified mirror rather than a verified organizer-authored canonical release.  These boundaries limit public data access and the availability of qualitative examples.

% TRACE: S8-C02 | manuscript/PAPER_PLAN.md | C1--C6 boundaries; configs/labels.yaml | pragmatic labels; manuscript/OPEN_ISSUES.md | N-03
The main conclusions are also bounded by one adjudicated Vietnamese-language dataset with a documented social-media source component and one operational task formulation.  The six pragmatic labels define the measured setting, not a complete account of Vietnamese pragmatics or a guarantee of behavior in another domain, population, label scheme, or deployment context.  ViSoBERT is relevant Vietnamese social-media context but was not evaluated as a baseline.  This missing comparison should be treated as future evaluation, not as an implicit negative result.

% TRACE: S8-C03 | answer/results/multitask_ablation.json | $.seed; answer/results/low_resource_sarcasm.json | $.seed; answer/results/main_pragmatic.json | GPT runs; answer/results/calibration.json | confidence eligibility; answer/data_provenance/rationale_audit_waiver.json
Several experimental qualifications further limit interpretation.  The multi-task ablation is single-seed, so its pragmatic-versus-external-ordinary-task difference is an observed trade-off rather than a causal or robust effect.  Each low-resource sarcasm budget also has one registered seed and includes non-monotonic outcomes; it cannot establish data efficiency.  Prompted API baselines are single recorded runs rather than multi-seed training results.  Calibration is available only for systems with stored confidence scores.  Finally, the study has no manual rationale-faithfulness verification and no permission-safe qualitative error examples.  It therefore does not claim faithful rationales, broad retention or transfer, universal pragmatics coverage, or stable low-resource advantages.
